\documentclass[review,11pt,3p,authoryear,nopreprintline]{elsarticle}

\usepackage[T1]{fontenc}
\usepackage{amsmath,amssymb,amsthm}
\usepackage{mathrsfs}
\usepackage{bm}
\usepackage{xcolor}
\usepackage[colorlinks=true,citecolor=blue,linkcolor=blue,urlcolor=blue]{hyperref}

\newtheorem{theorem}{Theorem}
\newtheorem{proposition}{Proposition}
\newtheorem{corollary}{Corollary}
\newtheorem{lemma}{Lemma}
\newdefinition{example}{Example}
\newdefinition{definition}{Definition}
\newdefinition{remark}{Remark}
\newdefinition{property}{Property}
\newdefinition{condition}{Condition}

\begin{document}

\begin{frontmatter}
  \title{Failure Probability Bounds with Uncertain Surrogate Residual Distributions}

  \author[inst1]{Renlong Wang}
  \ead{rw106@rice.edu}
  \affiliation[inst1]{organization={Department of Civil and Environmental Engineering, Rice University},
    city={Houston},
    postcode={TX 77005},
    country={USA}}

  \begin{abstract}
  \end{abstract}

  \begin{keyword}
  \end{keyword}
\end{frontmatter}

\section{Introduction}

\section{Problem Description}
\label{sec-problem}

Let $X\sim\mathbb{P}$ be a random input taking values in $\mathcal{X}\subseteq\mathbb{R}^{n}$, and let $g:\mathcal{X}\to\mathbb{R}$ be a measurable limit-state function. The failure probability is defined as
\[ 
  p_f := \Pr\{g(X)\leq0\} = \mathbb{E}_{X\sim\mathbb{P}}\left[\mathbf{1}_{\{g(X)\leq0\}}\right] = \int_{\mathcal{X}}\mathbf{1}_{\{g(\bm{x})\leq0\}}\,\mathbb{P}(\mathrm{d}\bm{x}). 
\]
For fixed $g$, this functional is linear in $\mathbb{P}$.

Evaluating $g$ at an input $\bm{x}$ can be computationally expensive, particularly when each evaluation requires a high-fidelity simulation. A common approach is to replace $g$ with a computationally efficient surrogate model. Let
\[
  \mathcal{D}:=\{(\bm{x}_i,g(\bm{x}_i))\}_{i=1}^{m}
\]
denote the training dataset used to construct a measurable surrogate $\hat{g}_{\mathcal{D}}$. Throughout the following analysis, $\mathcal{D}$ and $\hat{g}_{\mathcal{D}}$ are treated as fixed.

We represent the surrogate approximation error by a scalar random variable $\xi$ with distribution $\mathbb{Q}\in\mathcal{Q}(\Xi)$, where $\Xi\subseteq\mathbb{R}$ is a Borel support set and $\mathcal{Q}(\Xi)$ denotes the set of probability measures supported on $\Xi$. The training residual realizations are
\[
  \xi_{\mathcal{D}}(\bm{x}_i):=g(\bm{x}_i)-\hat{g}_{\mathcal{D}}(\bm{x}_i),
  \qquad i=1,\ldots,m.
\]

Let $\Pi$ denote the joint probability distribution of $(X,\xi)$ on $\mathcal{X}\times\Xi$, with marginal distributions $\mathbb{P}$ and $\mathbb{Q}$. The corresponding surrogate-based failure probability is defined as
\begin{equation} 
    p_{\hat f} := \Pr_{\Pi}\left\{\hat{g}_{\mathcal{D}}(X)+\xi\leq0\right\}
    = \int_{\mathcal{X}\times\Xi}\mathbf{1}_{\{\hat{g}_{\mathcal{D}}(\bm{x})+\xi\leq0\}} \Pi(\mathrm{d}\bm{x},\mathrm{d}\xi). 
\end{equation}
For fixed $\hat{g}_{\mathcal{D}}$, this functional is linear in $\Pi$.

The surrogate-based failure probability depends on the assumed residual distribution. This dependence is particularly relevant for inputs whose surrogate predictions lie close to the failure threshold, where even small residual errors can change the inferred failure status. Uncertainty in the residual distribution can therefore translate into uncertainty in the estimated failure probability. We argue that, however, in complex engineering systems, the available information may be insufficient to specify the error distribution reliably. To address this limitation, we construct an ambiguity set $\mathscr{Q}$ of candidate error distributions intended to contain the unknown true error distribution $\mathbb{Q}^{\star}$ with high confidence. Each distribution in $\mathscr{Q}$ represents a plausible candidate for the unknown true error distribution. 

For fixed $\mathbb{P}$ and $\hat{g}_{\mathcal{D}}$, we seek the range of surrogate-based failure probabilities over admissible joint distributions whose error marginals belong to $\mathscr{Q}$. The upper and lower bounds are
\begin{equation}
  \overline{p}_{\hat f} := \sup_{\mathbb{Q}\in\mathscr{Q}} p_{\hat f},
  \label{eq-upper-bound}
\end{equation}
and
\begin{equation}
  \underline{p}_{\hat f} := \inf_{\mathbb{Q}\in\mathscr{Q}} p_{\hat f},
  \label{eq-lower-bound}
\end{equation}
respectively. The interval describes the range of failure probabilities over joint distributions with input marginal $\mathbb P$ and error marginal in $\mathscr Q$, i.e., if $\mathbb Q^\star\in\mathscr Q$, then $\underline{p}_{\hat f}\leq p_f\leq\overline{p}_{\hat f}$. 

Accordingly, the keys to our success are:
\begin{enumerate}
  \item \textbf{Ambiguity set design and statistical calibration.} Specify admissible error distributions and dependence models, and develop calibration procedures under explicit statistical conditions.
  \item \textbf{Tractable reformulations.} Derive exact and computationally tractable reformulations of the infinite-dimensional optimization problems that define the upper and lower bounds, without introducing approximation error.
\end{enumerate}


\section{Ambiguity Set Design}

We characterize uncertainty in the surrogate residual distribution
through an empirical cumulative distribution function (CDF).
Let
\[
  \mathcal C
  :=\{(\bm Z_j,g(\bm Z_j))\}_{j=1}^{N}
\]
be a calibration dataset, where
$\bm Z_1,\ldots,\bm Z_N\overset{\mathrm{i.i.d.}}{\sim}\mathbb P$
are independent of the training dataset $\mathcal D$.
The surrogate $\hat g_{\mathcal D}$ is fitted and fixed before
the calibration data are used. Define the signed calibration residuals
\[
  \xi_j^{\mathrm{cal}}
  :=g(\bm Z_j)-\hat g_{\mathcal D}(\bm Z_j),
  \qquad j=1,\ldots,N,
\]
and their empirical CDF
\[
  \widehat F_{N,\mathcal D}(t)
  :=\frac{1}{N}\sum_{j=1}^{N}
    \mathbf 1_{\{\xi_j^{\mathrm{cal}}\le t\}},
  \qquad t\in\mathbb R.
\]
Conditional on $\mathcal D$, these residuals are independent
and identically distributed according to
\[
  \mathbb Q_{\mathcal D}^{\star}
  :=\operatorname{Law}_{X\sim\mathbb P}
    \bigl(g(X)-\hat g_{\mathcal D}(X)\bigr).
\]
We construct an ambiguity set containing all residual distributions
whose CDFs lie within a prescribed uniform distance of
$\widehat F_{N,\mathcal D}$.

\begin{definition}[CDF-based ambiguity set]
  Given the empirical CDF $\widehat F_{N,\mathcal D}$
  and a radius $\varepsilon_N\ge0$, define
  \begin{equation}
    \mathscr Q
    :=
    \left\{
      \mathbb Q\in\mathcal Q(\Xi)
      \;\middle|\;
      \sup_{t\in\mathbb R}
      \left|
        F_{\mathbb Q}(t)-\widehat F_{N,\mathcal D}(t)
      \right|
      \le\varepsilon_N
    \right\},
    \label{eq-cdf-ambiguity-set}
  \end{equation}
  where $F_{\mathbb Q}(t):=\mathbb Q((-\infty,t])$.
  The Borel set $\Xi$ is assumed to contain the true residual
  support; we take $\Xi=\mathbb R$ when no smaller support
  can be justified.
\end{definition}

The radius $\varepsilon_N$ accounts for sampling uncertainty
in the empirical CDF. For a fixed calibration sample size
$N\ge1$ and a confidence parameter $\delta\in(0,1)$, choose
\[
  \varepsilon_N
  =
  \sqrt{\frac{\log(2/\delta)}{2N}}.
\]
The Dvoretzky--Kiefer--Wolfowitz--Massart inequality then yields
\[
  \Pr_{\mathcal C}\left\{
    \mathbb Q_{\mathcal D}^{\star}\in\mathscr Q
    \,\middle|\,\mathcal D
  \right\}
  \ge1-\delta.
\]
Thus, conditional on the fitted surrogate, the ambiguity set
contains its true residual distribution with probability
at least $1-\delta$ over the calibration data.

The first constraint limits the deviation of the error mean from $\mu_0$. The second bounds the expected squared deviation from $\mu_0$, accounting for both error variance and mean deviation. 

Before calibrating the size parameters of the moment-based ambiguity set, we state the following finite-sample guarantee, which does not relies on the specify form of ambiguity set.
\begin{theorem} 
  For $\delta\in(0,1)$, if the ambiguity set $\mathscr{Q}$ contains the true distribution $\mathbb{Q}^{\star}$ with probability at least $1 - \delta$, i.e.,
  \[
    \Pr \left\{ \mathbb{Q}^{\star} \in \mathscr{Q} \right\} \geq 1 - \delta,
  \]
  then
  \[
    \Pr \left\{ \underline{p}_{\hat{f}} \leq p_{\hat{f}}(\mathbb{Q}^{\star}) \leq \overline{p}_{\hat{f}} \right\} \geq 1 - \delta.
  \]
\end{theorem}

\begin{proof}
  By the definition of the lower and upper bounds, $\mathbb{Q}^{\star} \in \mathscr{Q}$ imples
  \[
    \underline{p}_{\hat{f}} \leq p_{\hat{f}}(\mathbb{Q}^{\star}) \leq \overline{p}_{\hat{f}}.
  \]
  Therefore,
  \[
    \Pr\left\{ \underline{p}_{\hat{f}} \leq p_{\hat{f}}(\mathbb{Q}^{\star}) \leq \overline{p}_{\hat{f}} \right\}
    \geq
    \Pr\left\{\mathbb{Q}^{\star}\in\mathscr{Q}\right\} \geq 1-\delta.
  \]
\end{proof}

The following proposition states conditions under which the true probability distribution of $\varepsilon$ falls in the ambiguity set $\mathscr{Q}$ by a proper setting of $\gamma_1$ and $\gamma_2$ for the fixed sample size $N$ and confidence level $1 - \delta$.

% \begin{proposition}
%   For a given $\delta$, let 
%   \begin{equation}
%     \gamma_1^{N} = \qquad \gamma_2^{N} = . 
%   \end{equation}
%   There then exists a positive integer $N_0$ depending on $\delta$ such that, for the true probability distribution $\mathbb{Q}^{\star}$ of surrogate residual $\varepsilon$,
%   \[
%     \Pr\left\{\mathbb{Q}^{\star}\in\mathscr{Q}(\gamma_1^{N},\gamma_2^{N})\right\} \geq 1-\delta
%   \]
%   for all $N \geq N_0$.
% \end{proposition}

\section{Tractable Reformulation}

\end{document}
